% ECONOMICS_PROBLEM: {"id": "G21-319", "title": "Digital deposit runs", "jel_code": "G21", "source_id": "OP-0322"}
% !TeX program = xelatex
\documentclass[11pt,a4paper]{article}
\usepackage[margin=23mm]{geometry}
\usepackage{fontspec}
\setmainfont{Latin Modern Roman}
\usepackage{amsmath,amssymb,mathtools}
\usepackage{xcolor,enumitem,booktabs,longtable,array,xurl,hyperref}
\definecolor{muted}{HTML}{53636C}
\hypersetup{colorlinks=true,urlcolor=blue,linkcolor=blue}
\setlength{\parindent}{0pt}
\setlength{\parskip}{5pt}
\newcommand{\R}{\mathbb R}
\newcommand{\E}{\mathbb E}
\newcommand{\N}{\mathbb N}
\newcommand{\ind}{\mathbf 1}
\newcommand{\Var}{\operatorname{Var}}
\newcommand{\Cov}{\operatorname{Cov}}
\newcommand{\supp}{\operatorname{supp}}
\DeclareMathOperator*{\argmin}{arg\,min}
\DeclareMathOperator*{\argmax}{arg\,max}
\newcommand{\entrymeta}[3]{{\sffamily\small\color{muted}\textbf{\texttt{#1}}\quad|\quad #2\par #3\par}}
\newcommand{\Problemref}[1]{\texttt{#1}}
\newcommand{\JELref}[2]{\texttt{#2} (JEL \texttt{#1})}
\begin{document}
\section*{Digital deposit runs}
\textbf{Problem ID:} \texttt{G21-319}\quad \textbf{JEL:} \texttt{G21}\par
% BEGIN ECONOMICS STATEMENT
\label{problem:OP-0322}
{\sffamily\small\color{muted}Original subject group: Banking, credit and financial stability\par}
\label{definitions:problem:30007604}
\textbf{Audit classification:} Published research agenda. The agenda explicitly asks about technology and depositor/creditor run dynamics. The withdrawal threshold, graph and communication interventions are editorial; strategic coordination variants require game theory.
\subsection*{Definitions and precise research target}
Fix banks \(b=1,\ldots,B\) and a depositor network with observed links or an explicitly modeled latent-network distribution. Each depositor observes signals about bank fundamentals and messages from contacts, then decides when to withdraw subject to a transaction delay \(d\). Let \(v\) denote the speed or reach of social communication and \(m\) an official communication policy. Define \(T_b(d,v,m)\) as the first time cumulative withdrawals exceed a specified fraction \(\rho\in(0,1)\) of initial deposits. The task is to identify the causal effects of lower transaction delays and wider communication on
\[\Pr(T_b\leq h\mid x),\]
for each horizon \(h\) and prespecified fundamental state \(x\), together with contagion to other banks. Separate information transmission from strategic coordination by specifying how signals and messages enter depositor beliefs. Account for endogenous bank pricing, liquid-asset holdings and insurance coverage. Establish exclusions for changes in digital access or communication exposure; uncontrolled simultaneous bad news cannot identify a technology effect. Evaluate which message policies reduce socially costly runs without concealing deteriorating fundamentals. Report heterogeneity across depositor concentration and insurance status, and bounds where the network or signal quality is unobserved.

\textbf{Additional formal definitions and scope (editorial).} Use either continuous time \(t\in[0,H]\) with measurable withdrawal processes or discrete dates and say which. Cumulative gross withdrawals \(A_b(t)\) are measured against fixed initial deposits \(D_b(0)>0\); \(T_b=\inf\{t:A_b(t)\geq\rho D_b(0)\}\), with \(\inf\varnothing=\infty\). Distinguish gross withdrawals, net deposit outflows and failure. A communication network is a graph or stochastic graph with a specified message arrival kernel; \(v\) changes that kernel and \(d\geq0\) changes the earliest execution date. Potential withdrawal paths include depositor and bank responses under a declared equilibrium selection. An official message policy maps observed histories to verifiable announcements or a specified noisy signal; policy feasibility includes the sender's information and any truthfulness restriction. Strategic coordination requires an explicit strategic model and is a game-theoretic variant.

For the identification part, declare an admissible class \(\Theta\) of structural models, including preferences, technologies, shock laws, measurement/sampling rules and any equilibrium selector. If \(O\) denotes the complete observable history and \(T(\theta)\) the target defined above, the identified set is \(\mathcal I_T=\{T(\theta):\theta\in\Theta,\ \mathcal L_\theta(O)=\mathcal L_{\rm obs}(O)\}\); point identification means this set is a singleton. A \(do(a)\) comparison replaces stated policy/technology equations while fixing the declared exogenous law and re-solving the remaining equations. These definitions do not assert that an identification theorem holds without further restrictions.
\subsection*{Variants and assumption changes}
\begin{enumerate}
\item \textbf{Information transmission (editorial).} Hold strategic withdrawal incentives fixed and vary message-arrival speed or transfer delays separately.
\item \textbf{Coordination and contagion (editorial; game-theoretic).} Add strategic deposit withdrawals across a graph, deposit insurance and official signals; bound run probabilities over admissible equilibrium selections.
\end{enumerate}
\subsection*{References and source audit}
\begin{enumerate}
\item Official January 9, 2026 web version; locator: Prudential Architecture / Resilience, depositor/creditor runs paragraph. Broad agenda; exact model editorial. URL: \url{https://www.bankofengland.co.uk/research/bank-of-england-agenda-for-research}.
\end{enumerate}
\begin{enumerate}\raggedright
\item\hypertarget{definitions:source:30007604:1}{} Bank of England. 2026. \emph{Bank of England Agenda for Research, 2025–2028}. Prudential Architecture / Resilience, depositor/creditor runs paragraph.
\url{https://www.bankofengland.co.uk/research/bank-of-england-agenda-for-research}.
\end{enumerate}

\subsection*{Common conventions: Source mathematical conventions}

These conventions apply to each entry unless explicitly replaced.

\textbf{Models and identification.} A primitive model \(m\) specifies all probability laws, feasible choices, information, equilibrium and measurement rules used by its target. The admissible class \(\mathcal M\) is fixed independently of outcomes. If \(P_O(m)\) is its observable probability law and \(\tau(m)\) the target, its identified set at observable law \(P\) is
\[
 \mathcal I_\tau(P)=\{\tau(m):m\in\mathcal M,\ P_O(m)=P\}.
\]
Point identification means that this set is a singleton. An empty set signals incompatibility of data and assumptions. Coordinate bounds need not describe the joint set. Bounds are sharp only if every retained target is attainable or arbitrarily approachable by compatible models. Finite bounds require explicit restrictions on unobserved quantities; unrestricted confounding, hidden wealth, extrapolation or arbitrary equilibrium selection can yield uninformative sets.

\textbf{Causal interventions.} A potential outcome \(Y_i(p)\) is the outcome under a complete policy regime \(p\), including timing, financing, information and market spillovers. The target population and units are fixed before treatment. With interference, use the complete assignment vector or a justified exposure mapping; individual no-interference notation alone cannot identify an economy-wide intervention. A difference of mean potential outcomes does not require the cross-world joint distribution. Conditioning on post-treatment migration, survival, enrollment or employment changes the estimand and needs separate justification.

\textbf{Statistical statements.} An estimator, confidence set, forecast or test specifies a sampling law, model class and loss/coverage criterion. Uniform \((1-\alpha)\)-coverage means \(\inf_{m\in\mathcal M}\Pr_m\{\tau(m)\in C_n\}\geq1-\alpha\), or an explicitly asymptotic version. The whole parameter space trivially has coverage, so informativeness requires a stated width, risk or power objective. A forecasting improvement is relative to a named benchmark, information set and evaluation population.

\textbf{Equilibrium and optimization.} An equilibrium comprises feasible plans, optimality, beliefs where needed, budget constraints and all market/resource clearing. The primitives or a stated benchmark define each correspondence \(\mathcal E(p;m)\). Nonemptiness is assumed or proved, never inferred just from compact policy sets. A selected-equilibrium target uses a declared selection rule; a robust target quantifies over all permitted equilibria. Use suprema or infima unless attainment is established. An \(\varepsilon\)-optimal policy is feasible and within \(\varepsilon\) of the finite optimum value. Report an empty feasible set separately.

\textbf{Welfare and accounting.} Normative utility, interpersonal normalization, social weights, numeraire, discounting and the use of public receipts are primitives. Behavioral utility need not coincide with the welfare criterion. Household welfare includes ownership income and rents once; adding profits again to compensating variation generally double-counts them. A monetary equivalent is defined by an explicit transfer and price convention, with monotonicity and range conditions. Average effects, mechanism contributions, transfers and welfare losses are different targets. Infinite sums require convergence and dynamic feasible plans require terminal or no-Ponzi conditions.

\textbf{Source versions.} A working-paper date, revision date and journal publication year are distinguished. A DOI for a journal version is not automatically the DOI of an earlier manuscript. Locators refer to the stated version and printed pages where specified. A metadata-only check does not verify a claim about a full-text conclusion. Every entry's source subsection narrows the provenance claim accordingly.
% END ECONOMICS STATEMENT
\end{document}
